\documentclass{article}

\usepackage[utf8]{inputenc}  % Para caracteres especiais
\usepackage[portuguese]{babel} % Para a língua portuguesa

\usepackage[a4paper,top=2cm,bottom=2cm,left=2cm,right=2cm,marginparwidth=1.75cm]{geometry}

% Pacotes úteis
\usepackage{amsmath}
\usepackage{graphicx}
\usepackage[colorlinks=true, allcolors=blue]{hyperref}

\title{Título do trabalho}
\author{Identificação do(a) autor(a)}

\begin{document}
\maketitle

% O Abstract não é necessário
% \begin{abstract}
% Your abstract.
% \end{abstract}

\section*{Uso de IA Generativa}

Escreva nesta seção um texto conforme uma das duas opções apresentadas a seguir:

\textit{Opção 1: caso não tenha usado ferramenta de IA Generativa para composição de parte ou todo o texto da proposta:}

Declaro que nenhuma ferramenta de inteligência artificial generativa foi utilizada na elaboração desta proposta.

\textit{Opção 2: caso tenha feito uso de ferramenta de IA Generativa para composição de parte ou todo o texto da proposta:}

Declaro que a(s) ferramenta(s) de inteligência artificial generativa [nome] foi/foram utilizada(s) para [descrever as partes do texto nas quais a(s) ferramenta(s) foi/foram utilizada(s), e com que propósito].

\section{Introdução}

A introdução deve descrever o problema/pergunta de pesquisa que se quer abordar. Deve-se explicar o contexto do problema, incluindo menção a como este problema e/ou problemas similares são atacados, tanto na literatura acadêmica quanto na prática, se for o caso. Também deve conter uma indicação de como se pretende abordar a questão, da forma mais concreta possível. Lembre-se que o objetivo de um mestrado nunca é desenvolver um produto, mesmo que ele seja parte do desenvolvimento.

\section{Fundamentação teórica / trabalhos relacionados}

Descreva aqui conceitos essenciais para a compreensão do trabalho, assim como trabalhos importantes e relevantes que tratem de problemas similares ou iguais ao que será abordado. Não descreva conceitos básicos amplamente usados na área de pesquisa. Também não é preciso detalhar trabalhos individuais, exceto se houver grande ligação entre o trabalho citado e o que se pretende desenvolver. Por outro lado, tendências gerais e conceitos mais específicos do problema tratado devem ser detalhados, sustentados por citações de referências relevantes. Este é um exemplo de citação: \cite{kurose2021}.

\section{Metodologia}

Explique em linhas gerais como se pretende abordar o problema ou questão de pesquisa. Diagramas, figuras e exemplos podem ser úteis.

\section{Resultados esperados}

Descreva quais são os resultados e avanços pretendidos, assim como o seu impacto no mundo real e na área de pesquisa. Evite generalidades e lugares-comuns, concentre-se naquilo que o trabalho proposto pode de fato entregar.

\bibliographystyle{apalike}
\bibliography{sample}

\end{document}
